\documentclass[11pt,a4paper]{article}
\usepackage[T1]{fontenc}\usepackage[utf8]{inputenc}\usepackage{lmodern}
\usepackage[margin=23mm,headheight=15pt]{geometry}
\usepackage{amsmath,amssymb,amsthm,mathtools,booktabs,longtable,array,microtype,xcolor}
\usepackage{hyperref,fancyhdr,listings,enumitem}
\definecolor{ink}{HTML}{163348}\definecolor{accent}{HTML}{147D83}
\hypersetup{colorlinks=true,linkcolor=ink,urlcolor=accent,citecolor=accent,pdfauthor={Sébastien Palcoux}}
\pagestyle{fancy}\fancyhf{}\fancyhead[L]{\sffamily\small FUSION RING CENSUS}
\fancyhead[R]{\sffamily\small Rank \Rank}\fancyfoot[L]{\sffamily\footnotesize Computational companion \textbullet\ September 2026}
\fancyfoot[R]{\sffamily\small\thepage}\renewcommand{\headrulewidth}{0.4pt}
\newtheorem{theorem}{Theorem}[section]\newtheorem{proposition}[theorem]{Proposition}
\newtheorem{lemma}[theorem]{Lemma}\newtheorem{corollary}[theorem]{Corollary}
\theoremstyle{definition}\newtheorem{definition}[theorem]{Definition}\newtheorem{remark}[theorem]{Remark}
\newcommand{\ZZ}{\mathbb Z}\newcommand{\CC}{\mathbb C}\newcommand{\Q}{\mathbb Q}
\newcommand{\NN}{\mathbb Z_{\geq0}}\newcommand{\FPdim}{\operatorname{FPdim}}
\setlength{\parindent}{0pt}\setlength{\parskip}{0.45em}\setlength{\emergencystretch}{3em}
\lstset{basicstyle=\ttfamily\footnotesize,breaklines=true,columns=fullflexible,frame=single,rulecolor=\color{accent},aboveskip=1em,belowskip=1em}
\setcounter{tocdepth}{2}
\newcommand{\Rank}{5}
\begin{document}
\thispagestyle{empty}{\sffamily\small\color{accent} EXACT ENUMERATION / DATA / VERIFICATION}\par\vspace{1cm}
{\sffamily\Huge\bfseries\color{ink} Rank-5 fusion rings\par}
\vspace{0.35cm}
{\sffamily\Large A reproducible census through multiplicity 19\par}
\vspace{0.65cm}{\large Sébastien Palcoux}\par{\small BIMSA}\par\vspace{0.35cm}{\small Computational companion, 20 September 2026}\par\vspace{0.8cm}
\begin{abstract}
We explain the exact enumeration of all based fusion rings of rank 5 through multiplicity 19, comprising 34,133 isomorphism classes. The account separates the mathematical coverage argument, the completed search, and independent verification of the emitted tensors. All duality types are included, and no categorifiability filter is applied. Code, machine-readable counts, full multiplication tensors and provenance records accompany the text.
\end{abstract}\vspace{0.4cm}\noindent\textbf{AI assistance and verification.} GPT-6 Astra Pro assisted with code, exposition and computational checks. Exhaustive-search arguments and exact independent tensor verification are recorded separately; no proof-assistant certification or peer-reviewed acceptance is claimed.\vfill\noindent\textbf{Scope.} A census of based rings, not a classification of tensor categories. Only completed whole-rank results are released.\newpage\tableofcontents\newpage

\section{Objects, equivalence and the counting convention}
A \emph{fusion ring} here is a free abelian group with a distinguished finite basis
$B=\{b_0,b_1,\ldots,b_{r-1}\}$, an associative multiplication with unit $b_0$,
and an involution $i\mapsto i^*$ fixing $0$. Its structure constants satisfy
\[
b_i b_j=\sum_k N_{ij}^{k}b_k,\qquad N_{ij}^{k}\in\NN,
\]
\begin{align}
N_{0j}^{k}=N_{j0}^{k}&=\delta_{jk},&N_{ij}^{0}&=\delta_{i^*,j},\\
N_{ij}^{k}&=N_{i^*k}^{j}=N_{kj^*}^{i}.&&\label{eq:reciprocity}
\end{align}
These conventions agree with the usual based-ring definition; see
\cite{EGNO,VS}. No assumption of categorifiability, commutativity or integrality
of Frobenius--Perron dimensions is imposed unless proved in the relevant rank.

A based isomorphism is a bijection of the distinguished bases preserving the
unit and multiplication. In tensor notation it is a permutation $q$ fixing $0$
such that
\[
 N_{ij}^{k}={N'}_{q(i)q(j)}^{q(k)}\quad\hbox{for every }i,j,k.
\]
The coefficient of the unit determines the involution, so such an isomorphism
automatically intertwines duality. The \emph{multiplicity} is the maximum entry
of the \emph{full} ordered multiplication tensor. It is at least one, including
for group rings. Put
\[
 c_r(m)=\#\{\text{rank-}r\text{ fusion rings of multiplicity exactly }m\}/\cong,
 \qquad C_r(M)=\sum_{m=1}^{M}c_r(m).
\]
A list through $M$ contains all rings of multiplicity at most $M$, not only
those of multiplicity equal to $M$. This distinction is essential when exporting
OEIS b-files.

\subsection{Why duality types can be searched separately}
An involution on the $r-1$ nonunit elements consists of $p$ disjoint transpositions
and $r-1-2p$ fixed points, with $0\leq p\leq\lfloor(r-1)/2\rfloor$. Fix one
representative $d_p$ of each type. An isomorphism between two tensors with
involution $d_p$ commutes with $d_p$. Hence it is sufficient to quotient that
stratum by the centralizer
\[
 Z_{S_{r-1}}(d_p),\qquad |Z_{S_{r-1}}(d_p)|=2^p p!(r-1-2p)!.
\]
Different $p$ cannot be isomorphic. Therefore a whole-rank census is complete
exactly when every required stratum has been searched exhaustively. Finishing
only the non-self-dual strata does not certify a whole-rank counting term.
\section{The certified release at rank 5}
The released bound is $M=19$ and the cumulative number of classes is $C_{5}(19)=34,133$. Every duality type is included. The accompanying full tensor file has one record per based-isomorphism class.


\begin{center}\begin{tabular}{rr}\toprule Exact multiplicity $m$ & $c_{\Rank}(m)$\\\midrule
1 & 16\\
2 & 37\\
3 & 82\\
4 & 134\\
5 & 209\\
6 & 336\\
7 & 477\\
8 & 733\\
9 & 863\\
10 & 1,082\\
11 & 1,383\\
12 & 1,948\\
13 & 2,211\\
14 & 2,554\\
15 & 2,936\\
16 & 4,157\\
17 & 4,255\\
18 & 5,080\\
19 & 5,640\\
\bottomrule\end{tabular}\end{center}

The complete-run evidence is available through \path{results/census.json} and the corresponding reports under \path{verification/}. Earlier ranks retain their original completed-run provenance; repository checks and newly executed runs are labelled separately.
% Derivation excerpt, retained from the rank-five census project.
% The current modular speedup is justified in completeness.md.
\section{An independent census}\label{sec:data}
\subsection{Commutativity and coordinates}
The following based-ring argument justifies the commutative specialization. Formal-codegree facts used here are recalled in \cite{Ostrik}.

\begin{lemma}\label{lem:commutativity}
Every rank-five fusion ring is commutative.
\end{lemma}
\begin{proof}
The complex algebra of a Frobenius based ring is
semisimple: the coefficient-of-unit functional gives the positive
inner product making the basis orthonormal and multiplication a
$*$-representation. A noncommutative algebra of dimension five with a
one-dimensional representation must be $\CC\oplus M_2(\CC)$.
Its unique one-dimensional character is Galois invariant, hence rational,
and its values are algebraic integers, hence integers.
This character is the Frobenius--Perron dimension character. Its values are positive. Consequently its formal codegree is an
integer $D\geq5$, the sum of five nonzero integer squares.
If $f$ denotes the formal codegree of the two-dimensional representation,
the standard identity $\sum_E\dim(E)/f_E=1$ gives
\[
 f=\frac{2D}{D-1}=2+\frac2{D-1}.
\]
Formal codegrees are algebraic integers. The displayed rational number
is not an integer for $D\geq5$, a contradiction.
\end{proof}

Fix the unit $0$ and one of the three dualities
\[
(0,1,2,3,4),\quad(0,1,2,4,3),\quad(0,2,1,4,3).
\]
Write $C_{ijk}=N_{ij}^{k^*}$. Commutativity and reciprocity say that
$C$ is a symmetric tensor, invariant under simultaneous duality.
The nonunit triples therefore give respectively $20,13,10$ variables,
ordered lexicographically by their smallest orbit representative.
All variables range independently in $[0,M]\cap\ZZ$ before
associativity is imposed. The unit and duality determine the other
entries. Commutativity of the five multiplication matrices is equivalent
to associativity: the common unit is cyclic for the span of the matrices,
and commuting matrices acting identically on the unit agree everywhere.
Direct expansion gives $21,12,13$ distinct quadratic equations,
respectively. The supplement regenerates these polynomials from the
axioms; no source census enters their construction.

\subsection{Self-dual rings: reconstruction from a Krylov matrix}
In the self-dual case use the ordered triples
\[
\begin{gathered}111,112,113,114,122,123,124,133,134,144,\\
222,223,224,233,234,244,333,334,344,444\end{gathered}
\]
as coordinates $v_0,\ldots,v_{19}$. By relabelling, arrange that
$v_1$ is the minimum of the twelve mixed repeated-index coefficients,
and $v_1\leq v_2\leq v_3$. The first nontrivial fusion matrix is
\[
N_1=\begin{pmatrix}0&1&0\\1&a&u^{\mathsf T}\\0&u&Q\end{pmatrix},\qquad
u=(v_1,v_2,v_3)^{\mathsf T},\quad
Q=\begin{pmatrix}v_4&v_5&v_6\\v_5&v_7&v_8\\v_6&v_8&v_9\end{pmatrix},
\quad a=v_0.
\]
For each of the nine bounded coordinates of $(u,Q)$ set
\[
K=[u,Qu,Q^2u],\quad s=u^{\mathsf T}u,\quad t=u^{\mathsf T}Qu,
\quad H=Q^2-aQ+uu^{\mathsf T}-I.
\]
For $i=1,2,3$, let $T_i$ be the bottom-right $3\times3$ block
of $N_{i+1}$; thus $(T_i)_{jk}=C_{i+1,j+1,k+1}$.
The vectors $e_i$ below are the three standard coordinate vectors.
Commuting with $N_1$ gives
\begin{align}
T_i u&=He_i,\label{eq:krylov}\\
T_i Qu&=(QH-sQ+u(Qu)^{\mathsf T})e_i,\nonumber\\
T_i Q^2u&=(Q^2H-sQ^2-tQ+(Qu)(Qu)^{\mathsf T}
                       +u(Q^2u)^{\mathsf T})e_i.\nonumber
\end{align}
These follow successively from the lower blocks of the commutator.
If $\det K\ne0$, multiplication by $\operatorname{adj}(K)$ reconstructs
all ten remaining coefficients, affinely in $a$. One integrality
condition is a linear congruence $A+aB\equiv0\pmod{\det K}$.
A gcd calculation gives every bounded solution for $a$, including the
case $B=0$. For each such solution all remaining divisibilities,
coefficient bounds and all $21$ associativity equations are checked.
Thus there is no search over the last ten coefficients in this stratum.

\subsection{Every singular stratum}
The singular cases are essential to completeness. Multiplication by a
ring element is zero if it annihilates the unit. Hence any polynomial
annihilating the unit under $N_1$ annihilates $N_1$ itself. Apply this to
the cyclic block generated by the unit.

If $\operatorname{rank}K=2$, write $Q^2u=\alpha u+\beta Qu$ and
$\lambda=\operatorname{tr}Q-\beta$. The invariant rational subspaces of
$Q$ show that $\alpha,\beta,\lambda$ are integers. The remaining
one-dimensional eigenspace of $Q$ has eigenvalue $\lambda$. The cyclic
block condition is
\begin{equation}\label{eq:ranktwo}
(\lambda^2-a\lambda-1)(\lambda^2-\beta\lambda-\alpha)
-\lambda(s\lambda+t-\beta s)=0.
\end{equation}
This is linear in $a$. All zero-coefficient cases are retained.

If $\operatorname{rank}K=1$, write $Qu=\lambda u$ and let
$x^2-Tx+U$ be the characteristic polynomial on $u^\perp$.
The cyclic block polynomial is
\[
p(x)=x^3-(a+\lambda)x^2+(a\lambda-1-s)x+\lambda.
\]
For two distinct complementary eigenvalues require $x^2-Tx+U\mid p$;
for a repeated eigenvalue $\eta=T/2$ require $p(\eta)=0$.
These are again linear conditions on $a$, with no division by a
potentially zero coefficient.

For both positive-rank singular cases, the $15$ commutator equations
linear in the ten remaining coordinates are solved by fraction-free
elimination. Enumerate every bounded integral value of the free
coordinates, back-substitute, and test all $21$ quadratic equations.

If $u=0$, the unit is annihilated by $N_1^2-aN_1-I$.
For $a>0$, $x^2-ax-1$ is irreducible over $\Q$, and a five-dimensional
rational vector space cannot be a module over the resulting quadratic
field. Thus $a=0$ and $Q^2=I$. A nonnegative integral symmetric
orthogonal matrix is a permutation matrix. The transpositions are
handled by the same linear completion. For $Q=I$, the remaining ten
variables form the weighted rank-four associativity system: the constant
$1$ in each diagonal rank-four equation is replaced by $2$, because
both $1$ and the invertible element contribute. The Diophantine
elimination of the accompanying rank-four derivation applies with this substitution.
All zero branches are implemented explicitly in \texttt{weighted.hpp}.
For example, its terminal branch $b=c=e=0$ gives $a=1$, $d=f=2$ and
$gi+hj=h^2+i^2-6$; it must not be omitted.

\subsection{Nonself-dual rings and isomorphisms}
For one dual pair, use the $13$ orbit coordinates $v_0,\ldots,v_{12}$.
Put $A=v_6-v_5$, $B=v_{10}-v_9$, $C=v_{11}-v_{12}$.
Three associativity consequences are
\begin{align*}
Cv_5&=Av_2+Bv_4,& Cv_9&=Av_4+Bv_8,\\
v_{11}^2&=1+v_{12}^2-v_5^2+v_6^2-v_9^2+v_{10}^2.
\end{align*}
Enumerate $v_5,v_6,v_9,v_{10},v_{12}$ and recover $v_{11}$ by an exact
square test. Enumerate $v_4$ and solve for $v_2,v_8$, retaining all free
choices if $A$ or $B$ vanishes. Linear commutators determine or bound
the four remaining variables $v_0,v_1,v_3,v_7$; check all equations.
For two dual pairs the corresponding initial reductions are
\[
v_0^2=1+v_1^2-v_2^2-v_3^2+2v_4^2,\qquad
v_7^2=v_2^2-v_3^2+v_5^2.
\]
Two linear equations determine $v_6$ unless their coefficients vanish;
that branch retains every bounded value. Solve for $v_8,v_9$ by linear
completion and check all $13$ equations. The orbit definitions in
\texttt{census\_system2.json} and \texttt{census\_system4.json} fix these
coordinate conventions exactly.

Finally minimize over the permutations preserving the fixed duality:
$24,4,8$ permutations in the three cases. Distinct duality types cannot
be isomorphic. The independent verifier additionally minimizes every
output tensor over all $24$ unit-fixing permutations.

\subsection{The modular prefilter does not change the solution set}
In the nonsingular branch each reconstructed coefficient has the form
$(A+Ba)/D$. When $D$ is invertible modulo $257$, an allowed integer coefficient
must have a residue in $[0,M]$. Intersecting precomputed masks for two such
coefficients eliminates impossible $a$ values. Every retained value still
undergoes exact divisibility and the original integer checks. If $257\mid D$,
the original gcd/congruence calculation is used. All integer-singular branches
are unchanged. This is a necessary-condition speedup, not an assumption that
a generic determinant is nonzero.

\section{What the independent verification establishes}
The enumerator and verifier have separate mathematical tasks. Exhaustiveness
uses the complete coordinate parameterization, the safe-pruning proof and the
normal completion of every required branch. The verifier reads only the emitted
\emph{full tensors}; it neither reconstructs the parameterization nor trusts the
generator's polynomial list.

For every tensor it checks nonnegative integral coefficients, both unit laws,
uniqueness and involutivity of the dual, both reciprocity identities, and every
integer associativity identity for all four indices. It then canonicalizes
under \emph{every} unit-fixing basis permutation. A matching hash is only a
lookup aid: actual canonical tensors are compared. This detects both literal
repetitions and different labellings of an isomorphic based ring.

Passing these checks proves soundness and uniqueness of the emitted data; by
itself it would not prove that no ring was omitted. That is why normal completion
logs and the mathematical coverage argument are retained separately. Comparison
with the original ancillary file is an additional containment/regression test,
not an input to enumeration and not a replacement for an exhaustiveness proof.

\subsection{Provenance and reproducibility}
The retained Vercleyen--Slingerland v4 ancillary file has SHA-256
\begin{center}\footnotesize\ttfamily
b587c3f683157369da7cc090f762bbff0e18dcf11acc3f5031d6bc0c643843f8
\end{center}
Its 28,451 records contain 25,138 distinct classes. The 3,313 redundant records
occur at rank five and exact multiplicities 9--12. Their original and corrected
counts are respectively
\[
(1463,1794,2283,3049)\quad\longrightarrow\quad(863,1082,1383,1948).
\]
Every deletion has an explicit basis-permutation witness in the audit folder.
This audit does not establish completeness of the source's partial-search cells.

The primary manifest \path{results/census.json} records every released count,
bound, filename and checksum. Compressed tensors use the original line-per-ring,
nested-curly-brace format. An entry is $N[i][j][k]=N_{ij}^{k}$: the output index
has already been dualized where the internal parameter convention requires it.
The raw tensor needs no additional index conversion.

\subsection{Computation budgets and release gates}
The private repository provides manually dispatched GitHub Actions campaigns.
The original frontier has a 1,100-second shared driver deadline and a 20-minute
job ceiling. The long frontier selects only ranks 5--8: each rank has one
4,200-second (70-minute) shared deadline and a 75-minute job ceiling, at most
300 runner-minutes in total. Compilation, all attempted multiplicities and
duality types, export, compression and independent verification share the driver
deadline. Enumeration receives 90\% of remaining driver time; the balance is
reserved for other phases. Threads follow the actual runner CPU allocation.

Only a complete whole-rank run with successful independent verification enters
\path{best/}. A later incomplete attempt never replaces it. Candidates require
explicit review and fresh independent verification before integration. Partial
logs remain diagnostic evidence and never supply exhaustive census or OEIS terms.
All workflows are manual; neither a push nor this manuscript launches a search.

The reviewed long campaign (GitHub run 35440299307, revision
\texttt{931addc686e8}) completed rank five through 19: 34,133 classes, including
5,640 at exact multiplicity 19, with the old prefix unchanged. Larger attempts
at ranks five through eight timed out. The accepted result passed a fresh
independent GitHub audit before integration; see \path{docs/CAMPAIGN_REVIEW.md}
and \path{verification/imports/rank5_through19/} for source and audit records.

The later six-hour hosted rank-six job (35451839114) timed out. A laptop run
received on 20 September 2026 completed all duality types in 3,532.36 seconds
with 12 threads: 9,613 classes through 7, including $c_6(7)=3814$. The prefix
$(39,154,384,872,1582,2768)$ agrees. A fresh GitHub audit (35499152450) checked
all tensors and 1,153,560 unit-fixing permutations, finding zero duplicates.
See \path{docs/RANK6_LAPTOP_REVIEW.md} and
\path{verification/imports/rank6_through7/} for original logs and both audits.
Other ranks were not re-enumerated. Unlimited local runs require independent
verification and prefix agreement; GitHub workflow budgets remain finite.

\section{Reading and reusing the data}
The result files count based rings, not monoidal categories or realizations of
a ring. One tensor can admit several inequivalent categorifications, or none.
No number here is a classification of fusion categories. The rank-specific
count files use exact multiplicity; cumulative totals belong only in the
summary column. Proposed OEIS additions must be regenerated from a complete
manifest, not copied from an interrupted log.

The retained code was generated with AI assistance and is accompanied by explicit
mathematical coverage arguments and independent integer checks. These checks
are not a formal proof in a proof assistant. The present text is an explanatory
computational companion, not a claim of journal acceptance or external peer review.

\section{Reproduction commands}
From the repository root on a local computer, request a complete run with no clock limit. An interruption does not certify a smaller census.
\par\noindent\begin{minipage}{\linewidth}
\begin{lstlisting}
python3 scripts/run_laptop.py --rank 5 --bound 19
\end{lstlisting}
\end{minipage}\par
To verify the supplied release without recomputing it:
\par\noindent\begin{minipage}{\linewidth}
\begin{lstlisting}
python3 scripts/check_release.py --full
\end{lstlisting}
\end{minipage}\par
For a shared-budget frontier search, use \path{scripts/extend_census.py}. For ranks five through eight it starts at the released bound plus one and restarts the unfinished bound without claiming checkpoint resumption. The local command above verifies all tensors and the old count prefix before packaging a result. The separate timed frontier wrapper bounds every phase. Rank three is instead capped at 1000.
\clearpage

\begin{thebibliography}{9}\small\setlength{\itemsep}{2pt}
\bibitem{EGNO} P. Etingof, S. Gelaki, D. Nikshych and V. Ostrik,
\emph{Tensor Categories}, Mathematical Surveys and Monographs 205,
American Mathematical Society, 2015.
\bibitem{VS} G. Vercleyen and J. Slingerland,
\emph{On Low Rank Fusion Rings}, Journal of Mathematical Physics 64 (2023),
091703; \href{https://arxiv.org/abs/2205.15637v4}{arXiv:2205.15637v4}.
\bibitem{Ostrik} V. Ostrik, \emph{On formal codegrees of fusion categories},
Mathematical Research Letters 16 (2009), 895--901;
\href{https://arxiv.org/abs/0810.3242}{arXiv:0810.3242}.
\bibitem{DP} J. Dong and S. Palcoux, rank-four and rank-five census derivations,
project manuscripts and computational supplements, September 2026.
The relevant derivations are retained in the repository's \path{methods/} directory.
\bibitem{OEIS} OEIS Foundation Inc., \emph{The On-Line Encyclopedia of Integer
Sequences}, entries A348305, A354471, A354472, A354473, A354475, A354476
and A354477. Proposed updates in this repository are not submitted edits.
\end{thebibliography}
\end{document}
